\subsection{Equilibrio di Nash}
\label{subsec:EqNash}

Di seguito sono riportate alcune definizioni utili da \cite[Cap. 2]{MR4911797}. Si noti che l'analisi effettuata in questa sezione non cambia se si utilizzano le funzioni di costo come nell'\autoref{eq:10}.
\begin{definition}\label{def:1}
    Dato uno stato $\theta\in S^{\widehat{n}}\times S^{\widecheck{n}}$ chiamiamo \textit{hat}-rilevante, rispettivamente
    \textit{check}-rilevante, un tempo di viaggio $\widehat{T}_i(\theta)$, rispettivamente $\widecheck{T}_i(\theta)$, tale 
    che $\widehat{\theta}_i\neq 0$, rispettivamente $\widecheck{\theta}_i\neq 0$.
\end{definition}

\begin{definition}\label{def:2}
    Uno stato $\theta^{*}\in S^{\widehat{n}}\times S^{\widecheck{n}}$ è uno \textit{stato di equilibrio}
    se tutti i tempi di viaggio rilevanti, come da \autoref{def:1}, di ogni popolazione coincidono, ovvero se:
    \begin{equation*}
        \begin{split}
        \forall i,j\in\{1,...,\widehat{n}\} \text{, se }\widehat{\theta}_i^{*}\neq 0 \text{ e }\widehat{\theta}_j^{*}\neq 0
        \text{, allora }\widehat{T}_i(\theta^{*})=\widehat{T}_j(\theta^{*});\\
        \forall i,j\in\{1,...,\widecheck{n}\} \text{, se }\widecheck{\theta}_i^{*}\neq 0 \text{ e }\widecheck{\theta}_j^{*}\neq 0
        \text{, allora }\widecheck{T}_i(\theta^{*})=\widecheck{T}_j(\theta^{*}).
        \end{split}
    \end{equation*}
\end{definition}

La \autoref{def:2} assicura che all'equilibrio tutti i guidatori di una stessa popolazione impieghino lo stesso tempo per
arrivare dall'origine alla destinazione appartenenti a suddetta popolazione.

Grazie a queste definizioni ed a (\autoref{eq:3}), introduciamo:

\begin{definition}\label{def:3}
    Uno stato di equilibrio $\theta^{*}$ è un \textit{Equilibrio di Nash} se:
    \begin{equation*}
        \begin{split}
        \forall i\in\{1,...,\widehat{n}\}\text{ }\widehat{\theta}_i^{*}=0\implies\widehat{T}_i(\theta^*)\geq\widehat{T}_M(\theta^*);\\
        \forall i\in\{1,...,\widecheck{n}\}\text{ }\widecheck{\theta}_i^{*}=0\implies\widecheck{T}_i(\theta^*)\geq\widecheck{T}_M(\theta^*).
        \end{split}
    \end{equation*}
\end{definition}

\begin{definition}\label{def:4}
    È detto \textit{tempo di viaggio}, distintamente per le due popolazioni, il valore comune assunto dai $T$ rilevanti (\autoref{def:1})
    che caratterizzano l'Equilibrio di Nash trovato.
\end{definition}

Notiamo che, in uno stato di equilibrio $\theta^*$, $\widehat{T}_M(\theta^*)=\widehat{T}_i(\theta^*)$ $\forall i\in\{0,...,\widehat{n}\}$ (analogo per la popolazione \textit{check}),
dato che, per la \autoref{def:2}, tutti i tempi di viaggio dei percorsi sono uguali. Perciò la \autoref{def:3}
implica che un equilibrio di Nash è uno stato di equilibrio in cui a nessun guidatore di nessuna popolazione conviene cambiare percorso.
Ciò è dato dal fatto che in suddetto equilibrio i $T$ non rilevanti sono 
maggiori del tempo di viaggio, di conseguenza:
\begin{equation*}
    \begin{split}
    \forall i,j\in\{1,...,\widehat{n}\}\text{ }\widehat{\theta}_i^{*}=0,\text{ }\widehat{\theta}_j^*\neq0\implies\widehat{T}_i(\theta^*)\geq\widehat{T}_j(\theta^*);\\
    \forall i,j\in\{1,...,\widecheck{n}\}\text{ }\widecheck{\theta}_i^{*}=0,\text{ }\widecheck{\theta}_j^*\neq0\implies\widecheck{T}_i(\theta^*)\geq\widecheck{T}_j(\theta^*).
    \end{split}
\end{equation*}

Non è detto, però, che esista un equilibrio di Nash in ogni rete. Si riporta, dunque, da \cite[Cap. 2]{MR4911797}:
\begin{theorem}\label{th:1}
    Supponiamo che $\forall i\in\{1,...,N\}$ $(\widehat{\tau}_i,\widecheck{\tau}_i)\in C^0([0,1]^2;(\mathbb{R}_+\cup\{+\infty\})^2)$.
    Allora esiste un equilibrio di Nash $\theta^*\in S^{\widehat{n}}\times S^{\widecheck{n}}$, al senso della \autoref{def:3}.
\end{theorem}

\subsection{Funzioni ausiliarie}
\label{subsec:FunzAus}

Introduciamo le funzioni:
\begin{equation}
    \begin{split}
    \begin{array}{l c c c}
        \bar{\varphi}:\text{ } & \mathbb{R}_+\cup\{+\infty\} & \to & [0,1]\\
        & \xi & \mapsto &
        \begin{cases}
            \xi/(1-\xi)&\xi\in\mathbb{R}_+\\
            1 & \xi=+\infty
        \end{cases}
    \end{array}\\
    \begin{array}{l c c c}
        \varphi:\text{ } & (\mathbb{R}_+\cup\{+\infty\})^n & \to & [0,1]^n\\
        & (\xi_1,...,\xi_n) & \mapsto & (\bar{\varphi}(\xi_1),...,\bar{\varphi}(\xi_n)).
    \end{array}
    \end{split}
    \label{eq:4}
\end{equation}
Chiamiamo $C_+^n$ l'insieme dei $\theta\in\mathbb{R}$ tale che $\theta_i>0$ per almeno un $i\in\{1,...,n\}$.
Definiamo quindi la normalizzazione $\mathscr{N}:C_+^n\to S^n$, le cui componenti sono:
\begin{equation}
    \mathscr{N}(\theta)_i=\frac{\max\{0,\theta_i\}}{\sum_{j=1}^{n}\max\{0,\theta_j\}}\text{, per }i\in\{1,...,n\}
    \label{eq:5}
\end{equation}
Useremo $\varphi$ e $\bar{\varphi}$ indipendentemente che siano definite su $\mathbb{R}^{\widehat{n}}\cup\{+\infty\}$ o
su $\mathbb{R}^{\widecheck{n}}\cup\{+\infty\}$. Vale lo stesso per $\mathscr{N}$.

Nella dimostrazione del \autoref{th:1} presente in \cite[Cap. 4]{MR4911797} si fa utilizzo della funzione:
\begin{equation}
    \mathcal{F}:\text{ } S^{\widehat{n}}\times S^{\widecheck{n}} \to S^{\widehat{n}}\times S^{\widecheck{n}},
    \label{eq:6}
\end{equation}
che ha componenti:
\begin{equation}
    \begin{split}
    \begin{array}{l c c c}
        \widehat{\mathcal{F}}:\text{ } & S^{\widehat{n}}\times S^{\widecheck{n}} & \to & S^{\widehat{n}}\\
        & \theta & \mapsto & \mathscr{N}(\widehat{\theta} - \lambda((\varphi\circ\widehat{T})(\theta)-\widehat{\theta}^{\intercal}(\varphi\circ\widehat{T})(\theta)\mathds{1}_{\widehat{n}})) 
    \end{array}\\
    \begin{array}{l c c c}
        \widecheck{\mathcal{F}}:\text{ } & S^{\widehat{n}}\times S^{\widecheck{n}} & \to & S^{\widecheck{n}}\\
        & \theta & \mapsto & \mathscr{N}(\widecheck{\theta} - \lambda((\varphi\circ\widecheck{T})(\theta)-\widecheck{\theta}^{\intercal}(\varphi\circ\widecheck{T})(\theta)\mathds{1}_{\widecheck{n}})) 
    \end{array}
    \end{split}
    \label{eq:7}
\end{equation}
con:
\begin{equation}
    \lambda=\frac{1}{2}\min\left\{\frac{1}{\widehat{n}},\frac{1}{\widecheck{n}}\right\}.
    \label{eq:8}
\end{equation}

Per la dimostrazione del \autoref{th:1} ci si serve del Teorema del Punto Fisso di Brouwer, usato anche da John Nash nella sua tesi di
dottorato \cite[Cap. 3]{nash1950noncooperative} per dimostrare che \textit{"Every finite game has an equilibrium point"}, punto di equilibrio
denominato in seguito Equilibrio di Nash.

In \cite[Cap. 4]{MR4911797} viene dunque verificato che (\autoref{eq:6}), per Brouwer, ammette almeno un punto fisso:
\begin{equation}
    \exists\theta^*\in S^{\widehat{n}}\times S^{\widecheck{n}}:\text{ } \mathcal{F}(\theta^*)=\theta^*
    \label{eq:9}
\end{equation}
e che ogni punto fisso di $\mathcal{F}$ è un Equilibrio di Nash per la rete stradale, al senso della \autoref{def:3}.
